\section{Preliminaries }

In this section, we review existing methods for constructing reject-option predictors, illustrating them with a simple running example. We cover standard approaches, including maximum likelihood and Bayesian learning, and show how they apply to reject-option prediction. While the material is well established (e.g.,~\cite{Bishop-PRML-2006}), we present it using consistent notation and terminology to provide the necessary context for our novel framework for epistemic reject-option predictors, which builds upon these foundations.

\subsection{Bayes predictor}

Let $\SX$ denote the set of inputs and $\SY$ the set of targets. We consider that the input-output pairs $(x, y)$ are generated by a random process with distribution $p(x, y)$ defined over the product space $\SX \times \SY$. The goal is to construct a predictor $h\colon\SX\rightarrow\SY$, that gets an input  $x\in\SX$ and produces a prediction $h(x)\in\SY$ of the true target. Given a prediction loss $\ell \colon\SY\times\SY\rightarrow\Re_+$, the performance of a predictor $h$ is evaluated by its expected risk:
\[
    R(h) = \EE_{(x,y)\sim p(x,y)}\big [ \ell(y,h(x)) \big ]
    %\int_{x}\int_{y} p(x,y)\,\ell(y,h(x))
%\footnote{To keep the notation concise, we ommit the differential elements (e.g. $dx$, $dy$) and the domains (e.g. $\SX$ and $\SY$) which will be clear from the context. }
\:.
\]
%
The Bayes predictor $h^*$ is the theoretically optimal predictor that minimizes the expected risk $R(h)$, and it reads:
\begin{equation}
\label{equ:BayesPredictor}
   h^*(x) = \argmin_{\hat{y}\in\SY} \EE_{y\sim p(y|x)}\big [ \ell(y,\hat{y})\big ] \:.
   %\argmin_{\hat{y}\in\SY} \int_{y} p(y\mid x) \, \ell(y,\hat{y}) \:.
\end{equation}
% 
The Bayes predictor is not tied to any learning procedure.

\subsection{Aleatoric reject-option prediction}

The classical reject-option prediction framework~\cite{Chow-RejectOpt-TIT1970} allows a predictor to abstain from making a prediction on inputs $x \in \mathcal{X}$ that are likely to incur high loss. In this setting, the reject-option predictor $q: \mathcal{X} \rightarrow \mathcal{Y} \cup \{ \text{reject} \}$ outputs either a predicted target $q(x) \in \mathcal{Y}$ or abstains by returning $q(x) = \text{reject}$.

To accommodate the rejection option, the prediction loss $\ell \colon \mathcal{Y} \times \mathcal{Y} \rightarrow \mathbb{R}_+$ is extended to a reject-option loss $\ell^\epsilon: \mathcal{Y} \times (\mathcal{Y} \cup \{\text{reject}\}) \rightarrow \mathbb{R}_+$, incorporating a fixed rejection cost $\veps \in \mathbb{R}_+$, defined as:
%
\begin{align}
\label{equ:rejectOptLoss}
\ell^\veps(y,d) = 
\begin{cases}
\ell(y,d) & \text{if }\; d\in\mathcal{Y},\\
\veps & \text{if }\; d = \text{reject}\:.
\end{cases}
\end{align}
%
The performance of the reject-option predictor $q$ is then measured by its expected risk under the reject-option loss:
\begin{equation}
  \label{equ:ExpRiskRejectOption}
   R^\veps(q) = \EE_{(x,y)\sim p(x,y)}\big [ \ell^\epsilon(y,q(x)) \big ] \:.
   %\int_{x}\int_{y} p(x,y)\,\ell^\epsilon(q(x), y)\:.
\end{equation}
%
The optimal reject-option predictor $q^*$, which minimizes $R^\veps(q)$ is given by:
\begin{equation}
\label{equ:BayesOptimalPred}
   q^*(x) = \begin{cases}
   h^*(x) & \text{if }\; r^*(x)\leq \veps,\\
     \text{reject} & \text{if }\; r^*(x) > \veps\:,   
    \end{cases}
\end{equation}
%
where $h^*$ denotes the Bayes predictor~\equ{equ:BayesPredictor}, and 
\begin{equation}
\label{equ:condRisk}
   r^*(x) = \EE_{y\sim p(y|x)}[\ell(y,h^*(x))] \:,
   %\int_{\SY} p(y\mid x) \, \ell(y,h^*(x))\:,
\end{equation}
%
represents the conditional risk of $h^*$ at input $x$.

This formulation assumes complete knowledge of the data-generating distribution $p(x, y)$, and therefore accounts only for aleatoric uncertainty, while ignoring epistemic uncertainty. For this reason, we refer to the predictor in~\equ{equ:BayesOptimalPred} as the aleatoric reject-option predictor. In this context, the conditional risk $r^*(x)$ can serve as an optimal measure of aleatoric uncertainty in $x$ using the true model $p(x,y)$.

\begin{example}
\label{example:1}
Let $\mathcal{X} = \mathbb{R}$ be the set of inputs and $\mathcal{Y}=\Re$ the set of targets. Assume $(x, y) \in \mathcal{X} \times \mathcal{Y}$ is a sample from the distribution $p(x, y) = p(x)\, p(y \mid x)$, where
$p(x) = \mathcal{N}(x \mid 0,\, 1)$ is a standard normal distribution and $p(y \mid x) = \mathcal{N}(y \mid \mu(x),\ v(x))$ is normal with input-dependent mean $\mu(x)=0.5x + 1$ and variance $v(x) = 0.1 + 0.04(x + 8)^2$. 
Assume the squared loss $\ell(y, \hat{y}) = (y - \hat{y})^2$. Then, the Bayes predictor~\equ{equ:BayesPredictor} is given by the conditional mean:
$$
h^*(x) = \mathbb{E}_{y \sim p(y \mid x)}[y] = \mu(x)=0.5x + 1.
$$
Accordingly, the conditional risk~\equ{equ:condRisk} is equal to the conditional variance of $y$ given $x$:
$$
r^*(x) = \mathrm{Var}_{y \sim p(y \mid x)}[y] = v(x) = 0.1 + 0.04(x + 8)^2.
$$
Figure~\ref{fig:aleatoricRejectOption} illustrates the aleatoric reject-option predictor~\eqref{equ:BayesOptimalPred} using the expressions for $h^*(x)$ and $r^*(x)$ from this example. Note that the resulting predictor is independent of the marginal distribution $p(x)$; however, $p(x)$ influences learned predictors, as demonstrated in following examples.
\end{example}

\begin{comment}
\begin{figure}[h!]
    \centering
    % --- First Figure ---
    \begin{minipage}{0.8\linewidth}
        \centering
        \includegraphics[width=\linewidth]{figures/figure_1_aleatoric_uncertainty.pdf}
        \captionof{figure}{Lorem ipsum dolor sit amet, consectetur adipiscing elit, sed do eiusmod tempor incididunt ut labore et dolore magna aliqua. Ut enim ad minim veniam, quis nostrud exercitation.}
        \label{fig:aleatoric}
    \end{minipage}

    \vfill

    % --- Second Figure ---
    \begin{minipage}{0.8\linewidth}
        \centering
        \includegraphics[width=\linewidth]{figures/figure_2_epistemic_uncertainty.pdf}
        \captionof{figure}{Lorem ipsum dolor sit amet, consectetur adipiscing elit, sed do eiusmod tempor incididunt ut labore et dolore magna aliqua. Ut enim ad minim veniam, quis nostrud exercitation.}
        \label{fig:epistemic}
    \end{minipage}

    \vfill

    % --- Third Figure ---
    \begin{minipage}{0.8\linewidth}
        \centering
        \includegraphics[width=\linewidth]{figures/figure_3_total_uncertainty.pdf}
        \captionof{figure}{Lorem ipsum dolor sit amet, consectetur adipiscing elit, sed do eiusmod tempor incididunt ut labore et dolore magna aliqua. Ut enim ad minim veniam, quis nostrud exercitation.}
        \label{fig:total}
    \end{minipage}
\end{figure}
\end{comment}

\begin{figure}[ht!]
    \centering
    % --- Left Column (contains the single figure on the left) ---
    % The [c] option vertically centers this minipage relative to the right one.
    \begin{minipage}[l]{\linewidth}
        \centering
        % This is FIGURE 1
        \includegraphics[width=\linewidth]{figures/figure_1_aleatoric_uncertainty.pdf}
        \captionof{figure}{The aleatoric reject-option predictor outputs the prediction $h^*(x)$ in regions where the conditional risk $r^*(x)$ does not exceed the rejection cost $\varepsilon$, and abstains otherwise. This example considers the squared loss $\ell(y, \hat{y}) = (y - \hat{y})^2$, the reject cost $\veps=1$ and full knowledge of data distribution $p(x, y)$ defined in Example~\ref{example:1}. Under this loss, the predictor $h^*(x)$ returns the conditional mean $\mu(x)$, and the conditional risk $r^*(x)$ equals the conditional variance $v(x)$, which is the optimal measure of aleatoric uncertainty.
        }
        \label{fig:aleatoricRejectOption}
    \end{minipage}
\end{figure}

\begin{figure*}[ht!]
    \begin{subfigure}[b]{0.47\linewidth}        
      \includegraphics[width=\linewidth]{figures/figure_3_total_uncertainty.pdf}
      \caption{Bayesian reject-option predictor}
      \label{fig:BayesianRejectOpt}
    \end{subfigure}
    \hfill
    \begin{subfigure}[b]{0.47\linewidth}
       \includegraphics[width=\linewidth]{figures/figure_2_epistemic_uncertainty.pdf}
       \caption{Epistemic reject-option predictor}
      \label{fig:EpistemicRejectOpt}
    \end{subfigure}
    %
    \caption{Fig~\ref{fig:BayesianRejectOpt} illustrates the Bayesian reject-option predictor~\equ{equ:TotalRejectBL}, which outputs the prediction $H_B(x, D)$ in regions where the total uncertainty $T(x, D)$ does not exceed the rejection cost $\varepsilon = 2$, and abstains otherwise. Fig~\ref{fig:EpistemicRejectOpt} shows the epistemic reject-option predictor~\equ{equ:OptEpistRejOptPred}, which outputs the prediction $H_B(x,D)$ in regions where the epistemic uncertainty $E(x,D)$ does not exceed the rejection cost $\delta=1$, and abstains elsewhere. This example considers the setting of Example~\ref{example:BayesianRejectOption}, with training data $D$ consisting of $m = 10$ examples drawn from $p(x,y)$ specified in Example~\ref{example:1}. Under this setup, the predictor is given by $H_B(x, D) = \#x^T \#\mu_m$, where $\#\mu_m$ denotes the mean of the predictive distribution, serving as an estimate of the true mean $\mu(x)$. The total uncertainty is computes as $T(x, D) = \#x^T \#\Sigma_m \#x + v(x)$, which decomposes into the sum of aleatoric and epistemic uncertainties: $T(x, D) = A(x, D) + E(x, D)$, where $A(x, D) = v(x)$ and $E(x, D) = \#x^T \#\Sigma_m \#x$. 
    }
    \label{fig:BayesEpistemicRejOpt}
\end{figure*}
            
\begin{comment}
\begin{figure}[ht!]
    \centering
    \begin{minipage}[r]{\linewidth}    
        % --- Top-Right Figure ---
        \begin{minipage}{\linewidth}
            \centering
            \includegraphics[width=\linewidth]{figures/figure_3_total_uncertainty.pdf}
            \captionof{figure}{The Bayesian reject-option predictor outputs the prediction $H_B(x, D)$ in regions where the total uncertainty $T(x, D)$ does not exceed the rejection cost $\varepsilon = 1$, and abstains otherwise. This example considers the setting of Example~\ref{example:BayesianRejectOption}, with training data $D$ consisting of $m = 10$ examples. The predictor is given by $H_B(x, D) = \#x^T \#\mu_m$, where $\#\mu_m$ is the mean of the predictive distribution, serving as an estimate of the true mean $\mu(x)$. The total uncertainty is $T(x, D) = \#x^T \#\Sigma_m \#x + v(x)$, which decomposes into the sum of aleatoric and epistemic uncertainties: $T(x, D) = A(x, D) + E(x, D)$, where $A(x, D) = v(x)$ and $E(x, D) = \#x^T \#\Sigma_m \#x$.}
            \label{fig:totalRejectOption}
        \end{minipage}
        \vfill        
        % --- Bottom-Right Figure ---
        \begin{minipage}{\linewidth}
            \centering
            \includegraphics[width=\linewidth]{figures/figure_2_epistemic_uncertainty.pdf}
            \captionof{figure}{The epistemic reject-option predictor outputs the prediction $H_B(x,D)$ in regions where the epistemic uncertainty $E(x,D)$ does not exceed the rejection cost $\delta$, and absteins otherwise. This examples considers the same setup as in Figure~\ref{fig:totalRejectOption}. The rejection cost $\delta=2$ 
            specifies the maximal acceptable conditinal regret. The epistemic uncertianty reads $E(x,D)=\#x^T\#\Sigma_m\#x$.}
            \label{fig:epistemicRejectOption}
        \end{minipage}
    \end{minipage}
\end{figure}
\end{comment}

\subsection{Non-bayesian learning}

The true data-generating distribution $p(x, y)$ is typically unknown. Therefore, predictors must be learned from training data $D = ((x_i, y_i) \in \mathcal{X} \times \mathcal{Y} \mid i = 1, \ldots, m)$, which are assumed to consist of a sequence of independent and identically distributed (i.i.d.) samples drawn from $p(x, y)$.

In Maximum-Likelihood (ML) learning, one first selects a parametric family of distributions $\mathcal{P} = \{ p(x, y \mid \theta) \mid \theta \in \Theta \}$, with the assumption that this family contains the true distribution. That is, there exists a parameter $\theta^* \in \Theta$ such that
$p(x, y) = p(x, y | \theta^*)$. Given the training data $D$, a ML parameter estimate $\hat{\theta} \in \Theta$ is obtained by maximizing the log-likelihood function:
$$
L(\theta) = \log p(D \mid \theta)\quad\text{where}\quad p(D\mid \theta) = \prod_{i=1}^m p(x_i, y_i \mid \theta).
$$
%
Let us define Bayes predictor for a known parameter~$\theta$:
\begin{equation}
    \label{equ:PluginBayesPredictor}
h(x, \theta) = \arg\min_{\hat{y} \in \mathcal{Y}} \mathbb{E}_{y \sim p(y \mid x, \theta)} \left[ \ell(y, \hat{y}) \right]\:.
\end{equation}
%
The result of learning is a plug-in Bayes predictor $\hat{h}(x) = h(x, \hat{\theta})$, using the ML estimate of the parameter $\hat{\theta}$ in place of the true $\theta^*$.
%
In settings where a reject option is incorporated, the plug-in conditional risk 
$\hat{r}(x) = r(x, \hat{\theta})$ is used to approximate $r^*(x)$, where
\begin{equation}
\label{equ:plugingCondRisk}
r(x, \theta) = \mathbb{E}_{y \sim p(y \mid x, \theta)} \left[ \ell(y, h(x, \theta)) \right]\:.
\end{equation}

%Note that the optimal predictor depends only on the posterior distribution $p(y \mid x)$. Therefore, it is common to assume that the input $x$ is independent of the model parameters $\theta$, i.e., $p(x, y \mid \theta) = p(y \mid x, \theta)\, p(x)$. Under this assumption, it suffices to learn the conditional posterior $p(y \mid x, \theta)$ to construct plug-in decision rules. In this setting, maximum likelihood estimation reduces to maximizing the conditional log-likelihood: $\log L(\theta) = \sum_{i=1}^m \log p(y_i \mid x_i, \theta)$. The standard approach to training neural networks, based on using the cross-entropy loss on top of the softmax output layer, is a direct instance of this ML framework.

%A typical example of this approach is training neural networks with cross-entropy loss applied to softmax-normalized class scores. The network models the posterior $p(y|x,\theta)$, and learning objective can be seen as the conditional log-likelihood $\log L(\theta)=\sum_{i=1}^m \log p(y_i|x_i,\theta)$. 

%Empirical Risk Minimization (ERM) offers an alternative to Maximum-Likelihood learning. It selects a hypothesis class $\mathcal{H} \subset \{ h: \mathcal{X} \to \mathcal{Y} \}$ and learns $\hat{h} \in \mathcal{H}$ by minimizing empirical risk $\hat{R}(h) = \frac{1}{m} \sum_{i=1}^m \ell(y_i, h(x_i))$ on training data $D$. With a reject option, ERM can either approximate the Bayes-optimal predictor $q^*$ directly or estimate the conditional risk $r^*$, which is then combined with a separately trained predictor $\hat{h}$ to form a reject-option predictor $\hat{q}$.

ML learning yields a point estimate of the predictor $\hat{h}$ (or $\hat{q}$) given the training data $D$. This approach accounts for aleatoric uncertainty and, under suitable conditions, guarantees that epistemic uncertainty vanishes as the number of training examples $m$ tends to infinity~\cite{Wackerly-2002}. However, it does not explicitly represent or quantify epistemic uncertainty, which becomes problematic when the training data is insufficient.% to ignore it.



%The maximum-likelihod learning produces a point estimate of the predictor $\hat{h}$ (or $\hat{q})$ given training data $D$. This approache allows to capture the aleatoric uncertainty, and under suitable conditions, it is guaranteed to reduce epistemic uncertainty to zero as the number of examples goes to infinity. However, they do not represent or quantify epistemic uncertainty, which becomes problematic when the training data is unsufficient to ignore it.

\begin{example}Assume $p(x,y|\#\theta)=p(x)\,p(y|x,\#\theta)$
and $p(y|x,\#\theta)=\SN\big (y|x\theta_1+\theta_0,v(x)\big )$, where $\#\theta=(\theta_0,\theta_1)\in\Theta=\Re^2$ are unknown parameters to be learned and $v(x)=0.1 + 0.04(x + 8)^2$ a known variance.  Given i.i.d. training data $D=((x_i,y_i)\in\Re^2|i=1,\ldots,m)$, the maximum-likelihood estimate of the parameters reads:
\[
   \hat{\#\theta}= \big (\#X^T\#\Sigma^{-1}\#X\big )^{-1} \#X\#\Sigma^{-1}\#y \:,
\]
%
where 
\begin{equation}
\label{equ:DataXy}
\#X=\left [ \begin{array}{cc}x_1& 1\\ \vdots&\vdots \\ x_m&1\\ \end{array}\right]\in\Re^{m\times 2}\,, \quad\#y=\left [ \begin{array}{c} y_1\\ \vdots \\ y_m \end{array}\right ]\in\Re^{m}\,, 
\end{equation}
%
represent the training data and 
\begin{equation}\label{equ:DataNoise}
\#\Sigma = {\rm diag}\Big( [v(x_1),\ldots,v(x_m)]\Big )\in\Re_{+}^{m\times m}
\end{equation}
%
the known heteroscedastic noise in the training inputs.

\end{example}


\subsection{Bayesian Learning}

Bayesian learning~\cite{Bishop-PRML-2006,Gelman-Bayesian-2013} is an established framework to design predictors from data that model both the aleatoric and the epistemic uncertainty. It relies on a parametric generative model in which the model parameters are treated as random variables. Let $\theta \in \Theta$ denote the parameter, $D \in \SD = (\SX \times \SY)^m$ the training data, $x \in \SX$ an input observation, and $y \in \SY$ its associated target. The full tuple $(\theta,D,x,y)$ is considered a sample from the joint distribution $p(\theta,D,x,y)$ defined over the product space $\Theta\times\SD\times\SX\times\SY$.

The goal in Bayesian learning is to construct a predictor $H \colon \SX \times (\SX \times \SY)^m \rightarrow \SY$ that, given a training set $D$ and a new input $x$, outputs a prediction $H(x, D) \in \SY$. The performance of such a predictor is evaluated using the Bayesian expected risk:
\begin{multline}
\label{equ:BayesLearnExpRisk}
R_B(H) = 
%\int_{\Theta} \int_{\SD} \int_{\SX} \int_{\SY} p(x, y, D, \theta) \, \ell(y, H(x, D)) \\
 \EE_{(\theta,D,x,y)\sim p(\theta,D,x,y)} \big[\ell(y, H(x, D))\big]\;.
\end{multline}
%
The optimal Bayesian predictor minimizes the Bayesian expected risk~\eqref{equ:BayesLearnExpRisk}. It is based on the predictive distribution
$p(y|x,D)$ and is defined as:
\begin{equation}
\label{equ:LernPredictorBayes}
H_B(x, D) = \argmin_{\hat{y} \in \SY} \EE_{y\sim p(y|x,D)}\big [ \ell(y, \hat{y}) \big ] \:.
%\int_{\SY} p(y \mid x, D) \, \ell(y, \hat{y}) \:.
\end{equation}
%
%Note that the Bayes predictor~\equ{equ:BayesPredictor} relies on the conditional distribution $p(y | x)$, whereas the Bayesian predictor~\equ{equ:LernPredictorBayes} uses $p(y | x, D)$, incorporating information from the observed data $D$.


%\begin{example}(Supervised scenario)
%\label{examples:SupervisedScenario}
%Assume that the observation $x$ is independent of the parameters $\theta$ and training data $D$ are i.i.d. samples. Under the these assumptions, the generative model $p(\theta,D,x,y)$ can be decomposed to:
%\begin{equation}
%  \label{equ:GenModelDecomposition}
%  p(\theta,D,x, y) = p(\theta)\,p(x)\, p(D \mid \theta)\,  p(y \mid x, \theta) \:.
%\end{equation}
%where 
%\[
%p(D \mid \theta) = \prod_{i=1}^m p(x_i)\,p(y_i \mid x_i, \theta).
%\]
%
%Under the decomposition~\equ{equ:GenModelDecomposition}, the Bayes-optimal learning predictor~\equ{equ:LernPredictorBayes} simplifies to:
%\begin{equation}
%   \label{equ:SL-BayesLearnPred}
%   H^*(x,D) = 
%   \argmin_{\hat{y}\in\SY} \EE_{\theta\sim p(\theta|D)}\EE_{y\sim p(y|x,\theta)} \big [\ell(y,\hat{y})\big] \:,  
%   \argmin_{\hat{y}\in\SY} \int_{\Theta}p(\theta|D) \int_{\SY} p(y|x,\theta) \, \ell(y,\hat{y}) 
%\end{equation}
%
%\end{example}

\begin{example}
\label{example:BLregression}Assume $p(x,y|\#\theta)=p(x)\,p(y|x,\#\theta)$
and $p(y|x,\#\theta)=\SN(y|x\theta_1+\theta_0,v(x))$, with unknown parameters $\#\theta=(\theta_0,\theta_1)\in\Re^2$, and known variance $v(x)=0.1+0.04(x+8)^2$. Let prior distribution over the parameters be 2D normal distribution $p(\#\theta)=\SN(\#\theta|\#0,\#\Lambda)$, where $\#\Lambda={\rm diag}([\lambda_1,\lambda_2])$ is a diagonal covariance matrix. Let the training data be $D=((x_i,y_i)\in\Re^2)|i=1,\ldots,m)$, drawn i.i.d.  
%
Under this setup, the predictive distribution is normal $p(y|x,D)=\SN\big(y|\#x^T\#\mu_m,\#x^T\#\Sigma_m\#x+v(x)\big)$, where $\#x=[x,1]^T$, and the posterior parameters are given by
\[
\#\Sigma_m=\big(\#X^T\#\Sigma^{-1}\#X+\#\Lambda^{-1} )^{-1}\,,\quad \#\mu_m=\#\Sigma_m\#X^T\#\Sigma^{-1}\#y \:.
\]
Here, $\#X$, $\#y$ and $\#\Sigma$ are defined in equations~\equ{equ:DataXy} and~\equ{equ:DataNoise}.
%, and $\#\Lambda={\rm diag([\lambda_1^{-2},\lambda_2^{-2}])}$. 
Assuming the squared loss $\ell(y,\hat{y})=(y-\hat{y})^2$, the Bayesian predictor corresponds to the mean of the predictive distribution $p(y|x,D)$:
\begin{equation}
\label{equ:BLpredLinReg}
   H_B(x,D) = \EE_{y\sim p(y|x,D)}[y]=\#x^T\#\mu_m \:.
\end{equation}
\end{example}

\subsection{Bayesian reject-option predictor}

Bayesian learning is flexible and well-suited for constructing reject-option predictors from data. Specifically, the goal is to construct a function $Q\colon\SX\times(\SX\times\SY)^m\rightarrow\SY\cup \{\rm reject\}$ that, given a training data $D\in\SD$ and an input $x\in\SX$, either produces a predicted target $Q(x,D)\in\SY$ or opts to reject the input by returning $Q(x,D)={\rm reject}$. 

The performance of the reject-option predictor $Q$ is assessed via the Bayesian expected risk under the reject loss $\ell^\veps(y,\hat{y})$ (see Equation~\equ{equ:rejectOptLoss}):
\begin{multline}
\label{equ:BayesLearnRejectRisk}
R_L^\veps(Q) = \EE_{(\theta,D,x,y)\sim p(\theta,D,x,y)}\big [\ell^{\veps}(y, Q(x, D)) \big ]\:.
%\int_{\theta} \int_{D} \int_{x} \int_{y} p(x, y, D, \theta) \, \ell^{\veps}(y, Q(x, D))\:.
\end{multline}
%
The optimal Bayesian reject-option predictor $Q_B$ minimizing $R_L^\veps(Q)$ reads:
\begin{equation}
\label{equ:TotalRejectBL}
    Q_B(x,D)=\begin{cases}
       H_B(x,D) & \text{if }\, T(x,D) \leq \veps \:,\\
       \text{reject}   & \text{if }\, T(x,D) > \veps \:,\\
    \end{cases}
\end{equation}
%
where $H_B(x,D)$ is the Bayesian predictor~\equ{equ:LernPredictorBayes}, and 
\begin{equation}
\label{equ:TotalCondRisk}
    T(x,D) = \EE_{y\sim p(y|x,D)}\big[\ell(y,H_B(x,D))\big ]\:,
    %\int_{y}p(y\mid x,D) \, \ell(y,H^*(x,D)) \:,
\end{equation}
%
is the conditional risk for a given input $x$ and training data $D$. As discussed below, $T(x,D)$ reflects both aleatoric and epistemic uncertainty, and~\equ{equ:TotalCondRisk} is therefore referred to as the total uncertainty.


%The Bayesian reject-option predictor $Q^*$ makes accept-or-reject decisions based on the Bayesian conditional risk $T^*$, which serves as a measure of  total predictive uncertainty given $x$ and $D$. This risk captures both aleatoric and epistemic components of uncertainty. This becomes evident when expressing the predictive uncertainty as 
%\[
%   p(y \mid x, D) = \mathbb{E}_{\theta \sim p(\theta \mid D)} \left[ p(y \mid x, \theta) \right],
%\]
%
%under the assumption that $x$ and $\theta$ are independent. Here, the predictive distribution is a weighted average of $p(y | x, \theta)$, which captures aleatoric uncertainty. The weights $p(\theta | D)$ reflect our belief about the parameters given the data, that is, the epistemic uncertainty. 

\begin{example}
\label{example:BayesianRejectOption}
Consider the setup from Example~\ref{example:BLregression}. The Bayesian reject-option predictor~\equ{equ:TotalRejectBL} consists of the Bayesian predictor $H_B(x, D)$ defined in~\equ{equ:BLpredLinReg}, and the total uncertainty~\equ{equ:TotalCondRisk}, which in this setting takes the form
\begin{equation}
\label{equ:TotalCondRiskLinReg}
  T(x, D) = \mathrm{Var}_{y \sim p(y \mid x, D)}[y] = \#x^T\, \#\Sigma_m\, \#x + v(x) \:.
\end{equation}
%
%This expression corresponds to the total variance of the predicted target.
The resulting Bayesian reject-option predictor is illustrated in 
Figure~\ref{fig:BayesianRejectOpt}.
\end{example}


%Note that the total variance is a sum of two components $\#x^T \#\Sigma_m\#x$ and $v(x)$. The latter, $v(x)$, is the variance of the distribution $p(y|\mu(x),v(x))$ and can be seen as aleatoric variance. Below we show that the former, $\#x^T \#\Sigma_m\#x$, is a natural measure of the epistemic uncertainty. 